\appendixitem{支配多项式映射}

设 V、W 是有限维向量空间。设 f 是从 V 到 W 的一个多项式映射（《代数》，第四章，§5，第 10 节，定义 4）。若 \(\psi \in A_{W}\)，则 \(\psi \circ f \in A_{V}\)（出处同前，命题 17）。映射 \(\psi \mapsto \psi \circ f\) 是从 \(A_{W}\) 到 \(A_{V}\) 的一个同态，称为与 f 相伴的同态。它的核由那些函数 \(\psi \in A_{W}\)（在 \(f(V)\) 上为零（因而在 \(f(V)\) 在 Zariski 拓扑中的闭包上也为零）者）组成。

定义 1. 多项式映射 \(f: \mathrm{V} \to \mathrm{W}\) 称为支配的，如果从 \(\mathrm{A_W}\) 到 \(\mathrm{A_V}\) 的与 \(f\) 相伴的同态是单射的。

鉴于前述，\(f\) 是支配的当且仅当 \(f(\mathrm{V})\) 在 Zariski 拓扑下在 W 中稠密。

命题 3. 设 k 代数闭。设 \(f : V \to W\) 是一个支配多项式映射。V 的任何稠密开子集在 f 下的像含有 W 的一个稠密开子集。

只需证明，对每个非零元素 \(\varphi\)（\(A_{V}\) 的元素），\(f(V_{\varphi})\) 含有 W 的一个稠密开子集。把 \(A_{W}\) 与 \(A_{V}\) 的一个子代数等同起来（借助与 \(f\) 相伴的同态）。存在非零元素 \(\psi\)（\(\mathrm{A_W}\) 的元素），使得每个同态 \(w: \mathrm{A_W} \to k\)（不零化 \(\psi\) 者）都扩张为同态 \(v: \mathrm{A_V} \to k\)（不零化 \(\varphi\) 者）（《交换代数》，第五章，§3，第 1 节，定理 1 的推论 3）。而现在这样的 \(w\)（相应地，\(v\)）可以等同于 \(\mathrm{W}_{\psi}\)（相应地，\(\mathrm{V}_{\varphi}\)）的一个元素，且说 \(v\) 是 \(w\) 的扩张就意味着 \(f(v) = w\)。因而 \(\mathrm{W}_{\psi} \subset f(\mathrm{V}_{\varphi})\)。证毕。

设 \(f: \mathrm{V} \to \mathrm{W}\) 是一个多项式映射，\(x_0 \in \mathrm{V}\)。映射 \(h \mapsto f(x_0 + h)\)（从 \(\mathrm{V}\) 到 \(\mathrm{W}\) 的映射）是多项式的。把它分解为齐次多项式映射的一个有限和：

\[
f (x _ {0} + h) = f (x _ {0}) + \mathrm{D} _ {1} (h) + \mathrm{D} _ {2} (h) + \dots
\]

其中 \(D_{i}: V \to W\) 是 i 次齐次的（《代数》，第四章，§5，第 10 节，命题 19）。线性映射 \(D_{1}\) 称为 f 在 \(x_{0}\) 处的切线性映射。我们把它记作 \(\mathrm{D}f(x_{0})\)。

命题 4. 设 \(f: \mathrm{V} \to \mathrm{W}\) 是一个多项式映射。设存在 \(x_0 \in \mathrm{V}\)，使得 \((\mathrm{D}f)(x_0)\) 是满射的。则 \(f\) 是支配的。

在 V 中和 W 中各施行一个平移，可以设 \(x_{0} = 0\) 且 \(f(x_{0}) = 0\)。于是 f 作为齐次元素之和的分解可写成

\[
f = f _ {1} + f _ {2} + \dots \text {with} \deg f _ {i} = i,
\]

且由假设线性映射 \(f_{1}\) 是满射的。设 f 不是支配的。则存在非零元素 \(\psi\)（\(A_{W}\) 的元素），使得 \(\psi \circ f = 0\)。设 \(\psi = \psi_{m} + \psi_{m+1} + \cdots\) 是 \(\psi\) 分解为齐次元素的分解，其中 \(\deg \psi_{i} = i\) 且 \(\psi_{m} \neq 0\)。则

\[
\begin{array}{l} 0 = \psi \circ f = \psi_ {m} \circ f + \psi_ {m + 1} \circ f + \dots \\ = \psi_ {m} \circ f _ {1} + \rho , \end{array}
\]

其中 \(\rho\) 是次数 \textgreater{} m 的齐次多项式映射的和。由此推出 \(\psi_{m} \circ f_{1} = 0\)。由于 \(f_{1}\) 满射，\(\psi_{m} = 0\)，矛盾。

推论. 若 k 代数闭且 f 满足命题 4 的假定，则 V 的任何稠密开子集在 f 下的像含有 W 的一个稠密开子集。

这由命题 3 与命题 4 得出。
% semantic-fragments: legacy-input-seam
\relax{}
